% !TeX program = xelatex
\documentclass[11pt,a4paper]{article}

\usepackage[UTF8,scheme=plain,fontset=fandol]{ctex}
\usepackage[margin=22mm]{geometry}
\usepackage{amsmath,amssymb}
\usepackage[hidelinks]{hyperref}

\DeclareMathOperator{\KL}{D_{\mathrm{KL}}}
\DeclareMathOperator{\JSD}{JSD}
\newcommand{\E}{\mathbb{E}}
\newcommand{\pdata}{p_{\mathrm{data}}}
\setlength{\parindent}{0pt}
\setlength{\parskip}{5pt}
\setlength{\emergencystretch}{2em}
\allowdisplaybreaks[1]

\begin{document}

\begin{center}
  {\Large\bfseries Assignment 3: Generative Adversarial Networks}\\[8pt]
  Name: 张逸林
  \hspace{8mm}
  Student ID: 2500013576
\end{center}

\section*{1.}

\paragraph{(a)}
\[
\min_G\max_D V(D,G)
=
\min_G\max_D
\left\{
\E_{x\sim\pdata}[\log D(x)]
+
\E_{z\sim p_z}[\log(1-D(G(z)))]
\right\}.
\]
The discriminator distinguishes real samples from generated ones; the
generator aims to match the real data distribution by fooling the discriminator.

\paragraph{(b)}
For fixed \(G\), let \(\mu\) be the common reference measure.
\[
V(D,G)=\int
\left[\pdata(x)\log D(x)+p_g(x)\log(1-D(x))\right]\,d\mu(x).
\]
Maximize pointwise with \(a=\pdata(x)\), \(b=p_g(x)\), and \(t=D(x)\).
For \(a,b>0\),
\[
h(t)=a\log t+b\log(1-t),\qquad
h'(t)=\frac{a}{t}-\frac{b}{1-t},\qquad
h''(t)=-\frac{a}{t^2}-\frac{b}{(1-t)^2}<0.
\]
The unique maximum satisfies
\[
h'(t)=0\iff a(1-t)=bt\iff t^*=\frac{a}{a+b}.
\]
For \(a>0,b=0\), \(t^*=1\); for \(a=0,b>0\), \(t^*=0\). Thus
\[
\boxed{D^*(x)=\frac{\pdata(x)}{\pdata(x)+p_g(x)}}
\qquad\text{whenever }\pdata(x)+p_g(x)>0.
\]
If \(a=b=0\), any \(D(x)\in[0,1]\) is optimal.

\paragraph{(c)}
Let \(s=\sigma(a)\). Since \(\sigma'(a)=s(1-s)\),
\[
\frac{d\ell_{\mathrm{sat}}}{da}
=-\frac{s(1-s)}{1-s}=-s,
\qquad
\frac{d\ell_{\mathrm{ns}}}{da}
=-\frac{s(1-s)}{s}=s-1.
\]
For fixed \(D,z\), let \(s=D(G_\theta(z))\). By the chain rule,
\[
\begin{aligned}
\nabla_\theta\ell_{\mathrm{sat}}
&=-s\,\nabla_\theta a(G_\theta(z)),\\
\nabla_\theta\ell_{\mathrm{ns}}
&=-(1-s)\,\nabla_\theta a(G_\theta(z)).
\end{aligned}
\]
When \(s\approx0\), the factor \(-s\approx0\) suppresses generator gradients,
whereas \(-(1-s)\approx-1\) avoids this suppression.

\clearpage
\section*{2.}

Write \(p=\pdata\), \(q=p_g\), \(m=(p+q)/2\), and
\(S=\{x:p(x)+q(x)>0\}\).
Substituting the optimal discriminator,
\[
\begin{aligned}
C(G)
&=V(D^*,G)\\
&=\int_S p(x)\log\frac{p(x)}{p(x)+q(x)}\,d\mu(x)
 +\int_S q(x)\log\frac{q(x)}{p(x)+q(x)}\,d\mu(x)\\
&=\int_S p(x)\left(\log\frac{p(x)}{m(x)}-\log2\right)\,d\mu(x)
 +\int_S q(x)\left(\log\frac{q(x)}{m(x)}-\log2\right)\,d\mu(x)\\
&=\KL(p\|m)+\KL(q\|m)-2\log2\\
&=\boxed{2\JSD(p\|q)-2\log2}.
\end{aligned}
\]
By nonnegativity of KL divergence,
\[
\begin{aligned}
C(G)&\ge-2\log2,\\
C(G)=-2\log2
&\iff \KL(p\|m)=\KL(q\|m)=0\\
&\iff p=m=q\quad\text{almost everywhere}\\
&\iff p=q\quad\text{almost everywhere}.
\end{aligned}
\]
Hence, over all generator distributions,
\[
\boxed{\min_G C(G)=-2\log2.}
\]

\section*{3.}

GigaGAN keeps adversarial training but extends the original
unconditional GAN to large-scale text-to-image synthesis.

\begin{enumerate}
\item \textbf{Text conditioning.}
It uses \(G(z,c)\) and \(D(x,c)\) instead of \(G(z)\) and \(D(x)\).
Pretrained CLIP text features guide style mapping and word-level
cross-attention.

\item \textbf{Generator architecture.}
The StyleGAN2-based generator combines modulated convolutions with
self- and cross-attention. Sample-adaptive kernel selection mixes a
filter bank using weights determined by the style code.

\item \textbf{Discriminator and losses.}
A multi-scale input/output discriminator supervises each image-pyramid
level independently. The non-saturating GAN loss is supplemented with
matching-aware losses for mismatched captions paired with real and
generated images, CLIP contrastive loss, and vision-aided adversarial
loss using pretrained CLIP features.

\item \textbf{High-resolution synthesis.}
Two separately trained, text-conditioned GANs perform \(64\times64\)
image generation and upsampling to \(512\times512\), respectively.
\end{enumerate}

\end{document}
